\section*{Chapter \ref{ch:b1:ocean-circulation} --- The Circulation of the Ocean}
% values: calc/ch29.py

\begin{solution}{exo:b1:ocean-circulation:1}
(a)~At \qty{4}{m/s} (first formula): $1000\,C_D = 0.29 + 0.775 + 0.481 = 1.546$,
$T = 1.3 \times \num{1.546e-3} \times 16 = \qty{0.0322}{Pa}$. At \qty{7}{m/s}
(second formula): $1000\,C_D = 0.60 + 0.497 = 1.097$, $T = \qty{0.0699}{Pa}$. At
\qty{15}{m/s}: $1000\,C_D = 1.665$, $T = \qty{0.487}{Pa}$. (b)~$0.487/0.0699 =
6.97$, against $(15/7)^2 = 4.59$: the drag coefficient itself grows with the wind,
because a stronger wind raises rougher waves that grip the air better.
\end{solution}

\begin{solution}{exo:b1:ocean-circulation:2}
(a)~Northern Hemisphere, \ang{90} to the right of a wind blowing east: towards the
south. (b)~The wind blows towards the south-west; to its right is the north-west.
(c)~Southern Hemisphere, \ang{90} to the left of a wind blowing east: towards the
north. (d)~The wind blows towards the north; to its left is the west, which is
offshore on a west coast. In case (d) the transport points offshore: the surface
water is replaced from below, and cold water upwells along the coast.
\end{solution}

\begin{solution}{exo:b1:ocean-circulation:3}
(a)~$31.7 + 3.5 - 17.9 = \qty{17.3}{\sverdrup}$, against the published
\qty{17.1}{\sverdrup}: the means of the three pieces add up to the mean overturning
within the rounding and the uncertainties ($\pm\qty{3.2}{\sverdrup}$).
(b)~Above \qty{1100}{m}: $16.6 + 0.4 = \qty{17.0}{\sverdrup}$ northward; below:
$12.0 + 5.9 = \qty{17.9}{\sverdrup}$ southward. The two nearly balance, as they must
when no water accumulates north of the section; the small difference is within the
uncertainties. (c)~$\num{17.1e6} \times \num{3.156e7} = \qty{5.40e14}{m^3}$, that is
\qty{5.40e5}{km^3} a year; $\num{1.3324e9}/(\num{5.40e5}) \approx 2470$ years.
\end{solution}

\begin{solution}{exo:b1:ocean-circulation:4}
$D_E = 7.6 \times 10/\sqrt{\sin\varphi}$ and $V_0 = 0.0127 \times
10/\sqrt{\sin\varphi}$. At \ang{15}: \qty{149}{m} and \qty{0.250}{m/s}
(\qty{2.5}{\percent} of the wind); at \ang{45}: \qty{90}{m} and \qty{0.151}{m/s}
(\qty{1.5}{\percent}); at \ang{60}: \qty{82}{m} and \qty{0.136}{m/s}
(\qty{1.4}{\percent}). Near the equator $\sin\varphi \to 0$ and both formulas
diverge; more fundamentally $f \to 0$, so that the Coriolis force needed to balance
the friction vanishes, and the time $1/f$ the balance needs to set in becomes longer
than the winds stay steady. Ekman restricted his formula to $|\varphi| \geq
\ang{10}$.
\end{solution}

\begin{solution}{exo:b1:ocean-circulation:5}
(a)~$f = 2\Omega\sin\ang{11} = \qty{2.783e-5}{s^{-1}}$; $M = 0.07/(\num{2.783e-5}) =
\qty{2515}{kg.m^{-1}.s^{-1}}$, \ang{90} to the right of a wind blowing west:
northward. (b)~$Q = (2515/1027) \times \num{5e6} = \qty{1.22e7}{m^3/s}$,
\qty{12.2}{\sverdrup}. (c)~The estimate lies within the uncertainty of the direct
measurement ($12.0 \pm 5.5$) and close to the value from long-term winds
($13.5 \pm 0.3$): the theorem accounts for the observed transport.
\end{solution}

\begin{solution}{exo:b1:ocean-circulation:6}
(a)~$f(\ang{2}) = \qty{5.09e-6}{s^{-1}}$; $M = 0.05/(\num{5.09e-6}) =
\qty{9824}{kg.m^{-1}.s^{-1}}$, \qty{9.57}{m^2/s}: at \ang{2} N towards the north
(right of a westward wind), at \ang{2} S towards the south (left). (b)~Over a strip
\qty{1}{m} long in \omterm{def:b1:spherical-astronomy:geographic}{longitude}, $2 \times 9.57 = \qty{19.1}{m^3/s}$ leave across the
two \omterm{def:b1:spherical-astronomy:geographic}{latitudes}, \qty{444.8}{km} apart: $w = 19.1/(\num{4.448e5}) = \qty{4.30e-5}{m/s}$,
about \qty{3.7}{m} a day. (c)~At the equator $f = 0$ and $|\vec M| = T/|f|$ is
infinite: the steady Ekman balance does not exist there, and the numbers above, very
sensitive to the \omterm{def:b1:spherical-astronomy:geographic}{latitude} chosen, are only an order of magnitude. What survives is
the sign: on either side the transport points away from the equator, so the surface
water diverges and cold water rises along it.
\end{solution}

\begin{solution}{exo:b1:ocean-circulation:7}
(a)~$f(\ang{15}) = \qty{3.775e-5}{s^{-1}}$: $M = 0.06/(\num{3.775e-5}) =
\qty{1590}{kg.m^{-1}.s^{-1}}$ towards the north; $f(\ang{40}) =
\qty{9.375e-5}{s^{-1}}$: $M = \qty{1280}{kg.m^{-1}.s^{-1}}$ towards the south.
(b)~Between the two \omterm{def:b1:spherical-astronomy:geographic}{latitudes}, \qty{2780}{km} apart, $\partial M_y/\partial y =
-(1590 + 1280)/(\num{2.78e6}) = \qty{-1.03e-3}{kg.m^{-2}.s^{-1}}$: by
\cref{prop:b1:ocean-circulation:pumping}, $w_E = \qty{-1.005e-6}{m/s}$, a downward
velocity of \qty{31.7}{m} a year. (c)~$\num{1.005e-6} \times \num{2.78e6} \times
\num{5e6} = \qty{1.40e7}{m^3/s}$, \qty{14.0}{\sverdrup} pumped down.
\end{solution}

\begin{solution}{exo:b1:ocean-circulation:8}
(a)~$H = 1027 \times 3992 \times \num{17e6} \times 12 = \qty{8.36e14}{W}$,
\qty{0.84}{PW}. (b)~It is \qty{70}{\percent} of the measured \qty{1.2}{PW}: the real
upper limb is warmer near the surface than the mean of \qty{15}{\celsius}, and the
\omterm{def:b1:ocean-circulation:gyre}{gyre} also carries heat. (c)~$1027 \times 3992 \times \num{17e6} \times 3 =
\qty{2.1e14}{W}$, \qty{0.21}{PW}, a quarter of (a). Only water that becomes cold
before it returns carries much heat: the heat transport of the Atlantic depends on
the formation of deep water.
\end{solution}

\begin{solution}{exo:b1:ocean-circulation:9}
Let $x$, $y$ and $1 - x - y$ be the fractions of AAIW, NADW and AABW. Salt and heat
balances, measured from AABW:
\[
  -0.52\,x + 0.33\,y = 34.731 - 34.67 = 0.061,\qquad
  2.9\,x + 3.6\,y = 2.78 - 0.4 = 2.38 .
\]
The determinant is $-0.52 \times 3.6 - 0.33 \times 2.9 = -2.829$, and Cramer's rule
gives $x = (0.061 \times 3.6 - 0.33 \times 2.38)/(-2.829) = 0.20$ and $y = (-0.52
\times 2.38 - 2.9 \times 0.061)/(-2.829) = 0.50$: the sample is \qty{20}{\percent}
AAIW, \qty{50}{\percent} NADW and \qty{30}{\percent} AABW, a point inside the
triangle of the three waters.
\end{solution}

\begin{solution}{exo:b1:ocean-circulation:10}
(a)~With $W = u + iv = V_0e^{i\pi/4}e^{(1+i)az}$, whose real and imaginary parts are
the formulas of the proposition, $W'' = (1+i)^2a^2W = 2ia^2W$. The two equations
combine into $A_z(u'' + iv'') = -fv + ifu = if(u + iv)$, that is $A_zW'' = ifW$,
which holds since $2ia^2A_z = if$ for $a^2 = f/(2A_z)$. (b)~$W'(0) =
(1+i)aV_0e^{i\pi/4} = \sqrt2\,aV_0\,e^{i\pi/2} = i\sqrt2\,aV_0$: the surface stress
$\rho A_zW'(0)$ points north, with the size $\rho A_z\sqrt2\,aV_0 = \rho A_z\sqrt2
\sqrt{f/(2A_z)}\;T/(\rho\sqrt{fA_z}) = T$. (c)~$\int_{-\infty}^0 W\,\mathrm dz =
V_0e^{i\pi/4}/((1+i)a) = V_0/(\sqrt2\,a)$, a real number: the transport has no
northward part, and its eastward part is $\rho V_0/(\sqrt2\,a) =
T/\sqrt{fA_z}\cdot\sqrt{A_z/f} = T/f$. This is the theorem: size $T/f$, \ang{90} to
the right of a wind blowing north.
\end{solution}

\begin{solution}{exo:b1:ocean-circulation:11}
(a)~A northward step $\mathrm dy = a_\oplus\,\mathrm d\varphi$ changes $f = 2\Omega
\sin\varphi$ by $2\Omega\cos\varphi\,\mathrm d\varphi$, so $\beta = 2\Omega
\cos\varphi/a_\oplus = 2 \times \num{7.292e-5} \times \cos\ang{30}/(\num{6.371e6}) =
\qty{1.98e-11}{m^{-1}.s^{-1}}$. (b)~$R/\beta = \num{e-6}/(\num{1.98e-11}) =
\qty{5.04e4}{m}$, about \qty{50}{km}. (c)~Interior: $\num{30e6}/(\num{5e6} \times
1000) = \qty{6.0e-3}{m/s}$, \qty{6}{mm/s}; boundary current: $\num{30e6}/(\num{5.04e4}
\times 1000) = \qty{0.60}{m/s}$, a hundred times faster, the ratio of the widths. The
Gulf Stream, at \qty{1.8}{m/s}, is faster still: the real current is narrower than
this estimate, its flow is concentrated near the surface, and its recirculation adds
to its transport.
\end{solution}

\begin{solution}{exo:b1:ocean-circulation:12}
(a)~Wherever the sea is deeper than \qty{1000}{m}, the column above \qty{1000}{m}
holds at most \qty{1}{km} of water, and elsewhere less: the volume above
\qty{1000}{m} is at most $\num{361.84e6} \times 1 = \qty{3.618e8}{km^3}$, so the
volume below is at least $\num{1.3324e9} - \num{3.618e8} = \qty{9.71e8}{km^3}$.
(b)~$\num{26e6} \times \num{3.156e7} = \qty{8.21e14}{m^3}$ a year,
\qty{8.21e5}{km^3}; $\num{9.71e8}/(\num{8.21e5}) \approx 1200$ years, the same order as
the thousand years of the conveyor. (c)~$(1/2)^{25/12.3} = 0.244$, about a quarter;
$\num{e-3} \times 25 \times \num{3.156e7} = \qty{7.9e5}{m}$, about \qty{790}{km}.
After 25 years the tritium has spread no farther than a slow drift would carry it,
while a quarter of it is left to be detected.
\end{solution}

\begin{solution}{pb:b1:ocean-circulation:1}
\begin{enumerate}
  \item A wind from the north blows towards the south, along the coast.
  \item $1000\,C_D = 0.60 + 0.071 \times 8 = 1.168$: $C_D = \num{1.168e-3}$.
  \item $T = 1.3 \times \num{1.168e-3} \times 64 = \qty{0.0972}{Pa}$.
  \item $C_D = \num{1.452e-3}$ and $T = \qty{0.272}{Pa}$: the stress grows by $2.80$,
        more than $(12/8)^2 = 2.25$, because $C_D$ itself increases with the wind.
  \item $\Omega = \qty{7.292e-5}{rad.s^{-1}}$; $f = 2\Omega\sin\ang{38} =
        \qty{8.979e-5}{s^{-1}}$.
  \item $2\pi/f = \qty{7.00e4}{s}$, \qty{19.4}{h}. The Coriolis force turns a
        current through a full circle in this time; before the wind has blown for
        about that long, the water oscillates instead of settling into the steady
        balance of the theorem.
  \item \ang{90} to the right of a wind blowing south: towards the west, the open
        sea.
  \item $M = T/f = 0.0972/(\num{8.979e-5}) = \qty{1082}{kg.m^{-1}.s^{-1}}$.
  \item $1082/1027 = \qty{1.054}{m^2/s}$ for each metre of coast.
  \item $D_E = 7.6 \times 8/\sqrt{\sin\ang{38}} = \qty{77.5}{m}$.
  \item $1.054/77.5 = \qty{0.0136}{m/s}$, \qty{1.36}{cm/s}.
  \item $V_0 = 0.0127 \times 8/\sqrt{\sin\ang{38}} = \qty{0.129}{m/s}$, \ang{45} to
        the right of the wind: towards the south-west. The drifters give about half
        of Ekman's value, some \qty{0.07}{m/s}.
  \item The water leaving the coast in the \omterm{def:b1:ocean-circulation:ekman}{Ekman layer} cannot be replaced across the
        coast, nor from above: it is replaced by water rising from below the layer
        near the coast, as in the cross-section drawn in
        \cref{sec:b1:ocean-circulation:upwelling}; the isotherms bend up towards the
        shore.
  \item $1.054 \times 1000 \times 86400 = \qty{9.11e7}{m^3}$ a day per kilometre of
        coast.
  \item $w = 1.054/(\num{2e4}) = \qty{5.27e-5}{m/s}$, \qty{4.55}{m} a day.
  \item $77.5/4.55 = 17$ days.
  \item The air is cooled from below by the cold water: the layer near the surface
        becomes cold and stable, the moisture may condense into fog and low stratus,
        and with no convection there is little rain on the coast.
  \item The upwelled water is rich in nutrients; they feed the phytoplankton, which
        feed the zooplankton and the fish.
  \item A wind blowing north drives a transport \ang{90} to its right, towards the
        east: surface water piles up against the coast and sinks (\omterm{def:b1:ocean-circulation:upwelling}{downwelling}); the
        \omterm{def:b1:ocean-circulation:upwelling}{upwelling} stops and warmer surface water returns to the coast.
  \item $f(\ang{20}) = \qty{4.988e-5}{s^{-1}}$: $M = \qty{1948}{kg.m^{-1}.s^{-1}}$,
        $\sin\ang{38}/\sin\ang{20} = 1.80$ times more.
  \item In the Southern Hemisphere the transport is \ang{90} to the left of the
        wind. To push it west, offshore, the wind must blow towards the north, from
        the south: in both cases the wind blows along the coast towards the equator.
  \item $1000\,C_D = 0.29 + 0.62 + 0.308 = 1.218$; $T = \qty{0.0396}{Pa}$; $M =
        \qty{441}{kg.m^{-1}.s^{-1}}$; the ratio is $0.407$ for both.
  \item Per metre of coast, $M c_p\,\Delta T = 1082 \times 3992 \times 6 =
        \qty{2.59e7}{W}$, \qty{25.9}{MW}; along \qty{1000}{km}, \qty{2.59e13}{W},
        \qty{25.9}{TW}.
  \item $\num{31.7e6}/1.054 = \qty{3.01e7}{m}$, about \qty{30000}{km}: \omterm{def:b1:ocean-circulation:upwelling}{coastal
        upwelling} moves far less water than a \omterm{def:b1:ocean-circulation:wbc}{western boundary current}.
  \item $Q = 1.054 \times \num{e6} = \qty{1.05e6}{m^3/s}$: the \omterm{def:b1:ocean-circulation:upwelling}{upwelling} transport
        along \qty{1000}{km} of coast is \qty{1.05}{\sverdrup}.
\end{enumerate}
\end{solution}
